\subsection{Backend}
\label{sec:backend}

Backend to serwer REST API napisany w TypeScripcie na platformie Node.js
z frameworkiem Express 5 \cite{nodejs,express}. Kod jest kompilowany do modułów
ECMAScript i uruchamiany w kontenerze opisanym w podrozdziale~\ref{sec:docker}.

\subsubsection{Warstwy aplikacji}

Kod podzielono na warstwy, z których każda korzysta tylko z warstwy leżącej
bezpośrednio pod nią:

\begin{itemize}
  \item \textbf{trasy} (\texttt{routes/}) przypisują metody i ścieżki HTTP
        do funkcji kontrolera;
  \item \textbf{kontrolery} (\texttt{controllers/}) odczytują dane z żądania,
        wywołują serwis i budują odpowiedź HTTP (kod statusu, treść JSON);
  \item \textbf{serwisy} (\texttt{services/}) zawierają logikę biznesową
        i nie znają protokołu HTTP;
  \item \textbf{repozytoria} (\texttt{repositories/}) jako jedyne komunikują się
        z bazą danych (Prisma) i magazynem plików (API S3).
\end{itemize}

Aplikację buduje funkcja \texttt{createApp(deps)}, która otrzymuje gotowe
zależności: konfigurację, logger, klienta bazy danych i klienta S3. Tworzy je
dopiero plik \texttt{server.ts}, który następnie uruchamia serwer HTTP. Dzięki
temu testy mogą zbudować aplikację z innymi zależnościami, np. z adresem
nieistniejącej bazy danych, i sprawdzić zachowanie systemu przy awarii
(podrozdział~\ref{sec:testy}).

\subsubsection{Walidacja konfiguracji}

Przy starcie zmienne środowiskowe są sprawdzane schematem biblioteki Zod
\cite{zod}. Sprawdzane są m.in.: format adresu bazy danych, poprawność adresu
URL magazynu MinIO, minimalna długość sekretu JWT (32 znaki) oraz to, czy
klucz główny szyfrowania jest zapisem base64 dokładnie 32 bajtów. Puste
wartości są traktowane jak brakujące. Przy błędnej konfiguracji proces wypisuje
listę \emph{wszystkich} nieprawidłowych zmiennych (bez ich wartości, żeby nie
ujawnić sekretów w logach) i kończy się kodem~1. Kontener z błędną
konfiguracją zatrzymuje się więc od razu i widać przyczynę, zamiast działać
i zgłaszać błędy przy pierwszym żądaniu.

\subsubsection{Obsługa błędów}

Wszystkie błędy trafiają do jednego middleware'u obsługi błędów, który zwraca
je w jednolitym formacie:

\begin{small}
\begin{verbatim}
{ "error": { "code": "VALIDATION_ERROR",
             "message": "Request validation failed",
             "details": [ { "path": "email", "message": "..." } ] } }
\end{verbatim}
\end{small}

Pole \texttt{code} jest stałym identyfikatorem, na podstawie którego frontend
wyświetla komunikat po polsku. Pole \texttt{message} jest opisem dla
programisty, a \texttt{details} zawiera np. błędy poszczególnych pól formularza.
Błędy przewidziane przez aplikację są klasami pochodnymi \texttt{AppError}
(tabela~\ref{tab:bledy}). Każdy inny wyjątek zamieniany jest na odpowiedź
500 \texttt{INTERNAL\_ERROR} z ogólnym komunikatem, a szczegóły (np. treść
błędu bazy danych) trafiają wyłącznie do logu. Express 5 przekazuje do tego
middleware'u również wyjątki z funkcji asynchronicznych, więc kontrolery nie
potrzebują bloków \texttt{try/catch}.

\begin{table}[htbp]
\centering
\caption{Kody błędów API}
\label{tab:bledy}
\small
\begin{tabular}{|l|c|p{0.5\textwidth}|}
\hline
\textbf{Kod} & \textbf{HTTP} & \textbf{Znaczenie} \\
\hline
\texttt{VALIDATION\_ERROR} & 400 & Niepoprawne dane wejściowe (\texttt{ValidationError}) \\
\hline
\texttt{INVALID\_JSON} & 400 & Treść żądania nie jest poprawnym JSON-em \\
\hline
\texttt{FORBIDDEN} & 403 & Brak uprawnień do operacji (\texttt{ForbiddenError}) \\
\hline
\texttt{NOT\_FOUND} & 404 & Nieznana ścieżka lub zasób (\texttt{NotFoundError}) \\
\hline
\texttt{PAYLOAD\_TOO\_LARGE} & 413 & Treść JSON większa niż 1~MB \\
\hline
\texttt{INTERNAL\_ERROR} & 500 & Nieoczekiwany błąd serwera \\
\hline
\end{tabular}
\end{table}

\subsubsection{Logowanie zdarzeń}

Logi są zapisywane biblioteką pino \cite{pino} jako jeden obiekt JSON na
linię. Taki format łatwo przeszukiwać i zbierać centralnie
(podrozdział~\ref{sec:monitoring}). Każde żądanie dostaje identyfikator
(\texttt{requestId}): jeśli klient przesłał poprawny nagłówek
\texttt{X-Request-Id}, jest on używany ponownie, w przeciwnym razie
generowany jest nowy UUID. Identyfikator trafia do nagłówka odpowiedzi i do
każdego wpisu w logu powstałego podczas obsługi żądania. Po zakończeniu
żądania zapisywany jest jeden wpis z metodą, adresem, kodem odpowiedzi i czasem
obsługi, na poziomie \texttt{warn} dla błędów klienta i \texttt{error} dla
błędów serwera. Nagłówki \texttt{Authorization} i \texttt{Cookie} są w logach
zamaskowane. Zapytania kontrolne Dockera do \texttt{/health}, wykonywane co
10~sekund, nie są logowane. W trybie deweloperskim logi są formatowane
czytelnie przez \texttt{pino-pretty}.

\subsubsection{Punkty kontroli stanu}

Backend udostępnia dwa punkty kontroli stanu (tabela~\ref{tab:endpointy}),
dostępne pod ścieżką \texttt{/health} (dla Dockera) oraz \texttt{/api/health}
(przez serwer nginx).

\begin{itemize}
  \item \textbf{Liveness} (\texttt{GET /health}) odpowiada kodem 200, dopóki
        proces obsługuje żądania HTTP. Celowo nie sprawdza bazy danych:
        awaria bazy nie powinna powodować restartu działającego backendu.
  \item \textbf{Readiness} (\texttt{GET /health/ready}) równolegle sprawdza
        PostgreSQL (zapytanie \texttt{SELECT 1}) i MinIO (\texttt{HeadBucket}
        na buckecie dokumentów), każde z limitem 3~sekund. Zwraca 200 i
        \texttt{\{"status":"ready"\}} albo 503 i \texttt{\{"status":"not\_ready"\}},
        a w polu \texttt{checks} status każdej zależności (\texttt{up} lub
        \texttt{down}). Przyczyny awarii są zapisywane w logu, ale nie są
        zwracane klientowi, żeby nie ujawniać wewnętrznych adresów.
\end{itemize}

\begin{table}[htbp]
\centering
\caption{Endpointy API}
\label{tab:endpointy}
\small
\begin{tabular}{|l|l|p{0.4\textwidth}|}
\hline
\textbf{Metoda} & \textbf{Ścieżka} & \textbf{Opis} \\
\hline
GET & \texttt{/health}, \texttt{/api/health} & Liveness -- proces działa \\
\hline
GET & \texttt{/health/ready}, \texttt{/api/health/ready} & Readiness -- baza danych i MinIO dostępne \\
\hline
\end{tabular}
\end{table}

\douzupelnienia{endpointy uwierzytelniania, profilu i dokumentów,
szyfrowanie kopertowe (diagram sekwencji uploadu), wersjonowanie, autoryzacja;
Etapy 4--6}
